% Auto-generated by make_appendix_tables.py — do not hand-edit.
\begin{table}[t]
\caption{Label-noise reference levels by surface and cloud-distance regime for the production anomaly targets: target variance and the maximum explainable $R^2$ under three noise scenarios (reference-mean sampling error only; additionally treating the retrieval posterior $\sigma^2$ as noise; treating all far-field variance as noise).}
\label{tab:label-noise}
\small
\begin{tabular}{lrrrrr}
\tophline
Subset & $n$ & var($y$) & $R^2_{\mathrm{max,ref}}$ & $R^2_{\mathrm{max,ref+ret}}$ & $R^2_{\mathrm{max,emp}}$ \\
\middlehline
\multicolumn{6}{l}{\textit{Ocean (r05 target)}} \\
all & 9\,271\,499 & 0.521 & 0.990 & 0.603 & 0.522 \\
near-cloud ($<$10\\,km) & 5\,558\,366 & 0.707 & 0.990 & 0.702 & -- \\
\middlehline
\multicolumn{6}{l}{\textit{Land (r15 target)}} \\
all & 4\,463\,809 & 0.663 & 0.993 & 0.606 & 0.460 \\
near-cloud ($<$10\\,km) & 512\,386 & 2.667 & 0.995 & 0.878 & -- \\
\bottomhline
\end{tabular}
\\[3pt]\parbox{\linewidth}{\footnotesize $R^2_{\mathrm{max,ref}}$ (reference-sampling noise only) is the only hard bound; the retrieval-posterior and empirical far-field scenarios are stated noise levels that the anomaly target can legitimately exceed (Sect.~3.4 and Fig.~C4). Computed on the pooled 140-date cohort (Table~\ref{tab:cohort-attrition}).}
\end{table}
